\begin{table}[!htbp]
\centering
\caption{\textbf{RE point benchmarks and documented competitive-price intervals in A/B}}
\label{tab:app-competitive-intervals}
\footnotesize
\fittable{\begin{tabular}{lllrrrrl}
\toprule
Market & State & Side & $\bar v$ & $v_{(2)}$ & $P_{\RE}=v_{(1)}$ & $\mathcal P^{CE}$ & $\mathcal D^{CE}$ \\
\midrule
A & X & Buying & 260 & 330 & 360 & $[330,360]$ & $[0.7,1.0]$ \\
A & Y & Selling & 260 & 130 & 160 & $[130,160]$ & $[1.0,1.3]$ \\
B & X & Buying & 300 & 380 & 400 & $[380,400]$ & $[0.8,1.0]$ \\
B & Y & Selling & 300 & 180 & 200 & $[180,200]$ & $[1.0,1.2]$ \\
\bottomrule
\end{tabular}}
\caption*{\footnotesize Notes: Valuations and price intervals are in francs per certificate. $v_{(2)}$ is the second-highest informed type valuation; $P_{\RE}$ is the highest and the upper price endpoint. $\mathcal D^{CE}$ maps the price interval through $D=(p-\bar v)/(P_{\RE}-\bar v)$. Source: the author-supplied note on the normalized axis and competitive-price benchmarks, Section 3.3, which records the market-parameter code convention. These intervals do not replace the paper's point-based outcome or classify observed transactions.}
\end{table}
